% \subsection{LOCA-Downscaled CMIP6 Precipitation Ensemble}


\subsubsection{Case Studies}
\paragraph{The Data}
To demonstrate that \ClimateSOM can uncover valuable insights, we apply the workflow to a ensemble dataset of LOCA-downscaled CMIP6 precipitation projections~\cite{pierce2014statistical,pierce2015improved,eyring2016overview} over California and the Northwest United States\cite{pierce2023future}.
%
% The Localized Constructed Analogs (LOCA) statistical downscaling method \cite{pierce2014statistical,pierce2015improved } is an analog-based method used to downscale
% relatively coarse scale global climate model projections from the CMIP6 archive \cite{eyring2016overview}  to finer scale regional projections. 
%
% The LOCA North American projections were carried out at 6km spatial resolution, in which downscaling of precipitation was conducted from global model projections made by 27 CMIP6 models, 
%
These model projections were made under 3 different SSPs \cite{riahi2017shared}, from a  SSP245, a relatively low greenhouse gas emission future, SSP370, a medium high greenhouse gas emissions future, and SSP585 a high greenhouse gas emissions future. 
%
% Of the 27 models, we consider \textasciitilde 14 models for each SSP and the historical period (i.e. four time periods per GCM) for which data was available, resulting in a total of 56 ensemble members.
We consider 14 GCMs for each SSP and the historical period (i.e. four time periods per GCM) for which data was available, resulting in a total of 56 ensemble members.
%
Details of the ensemble members used in this work can be found in the Appendix.
% LOCA or Localized Constructed Analogs, is a statistical downscaling method that takes coarse resolution climate model outputs, here of CMIP6\cite{eyring2016overview}, and downscales them to a regular grid of finer resolution.

We consider two regions separately: California and the Northwestern United States, encompassing Oregon, Washington, and Idaho.
%
While daily predictions are available, we focus on monthly projections considering months from October to May, as these months contribute to \textasciitilde85\% of the annual precipitation in both spatial regions.
%
We apply two preprocessing steps. 
%
Second, we normalize the monthly anomalies by dividing them by per-month standard deviation across models, time, and space.
%
The final data, therefore, consists of monthly precipitation anomalies expressed in standard deviations from the historical average.
%
For the California case study the each ensemble model run had a of size of $780\times 151\times 165$ (time $\times$ longitude $\times$ latitude), while for the Northwestern US, it was $780\times 60\times 115$.
%
We consult \textbf{E1}, a co-author, for the California case study, and another domain expert \textbf{E2} for the Northwestern US case study.


% While this makes any analysis across different months difficult, this simplification allows us to train a single SOM across different months that streamlines the analysis.

\paragraph{LOCA-Downscaled CMIP6 over California}
\paragraph{SOM Training}
We experimented with various dimensions of the SOM grid and found that the choice had minimal impact on the results so long they were sufficiently large and selected a 30$\times$30 grid.
%
We include the impacts of the initial neighborhood radius ratio $kR$ and the neighborhood decay ratio $kS$ on the SOM in the Appendix.
%
For the California region, we selected $kR=0.03$ and $kS=0.2$ as they provided a good balance explanatory power and SOM grid smoothness.
%
This provided a SOM grid with 80.5\% of explained variance.
\paragraph{The Case Study}
We outline a case study of the precipitation ensemble over California with the \textbf{E1}, a climate scientist and co-author of this paper.
%
We will reference~\cref{fig:Dan} for this case study.
%
First, the expert goes through \step{S1}: Anchor to reorganize the SOM node space to improve the interpretability of the SOM node space.
%
They place wetter precpitation SOM nodes to the left, dryer nodes to the right, and allow the vertical axis to represent the North-South gradient of the region, placing nodes with higher precipitation in the North on the upper side of the SOM Node Space and vice versa.
%
Next, as they are familiar with the SOM nodes and their precipitation patterns, they perform \step{S2}: Annotate producing five annotations regions: (1) \textit{Statewide Wet}, (2) \textit{Wetter North - Drier South}, (3) \textit{Wetter South - Drier North}, (4) \textit{Intermediate}, and (5) \textit{Statewide Dry} (\cref{fig:Step1and2}).
%
The first question the expert wishes to answer is on the behavior of model runs that predict a wetter California in January and whether there are trends that these model run exhibit in the previous months.
% 
To this end, the expert examines the model run clustering in the circuit-line diagram and identifies two clusters of model runs in January (\req{R3}), one representing a wetter California than the other based on the color of the cluster boxes.

They then compare the two January clusters (\req{R2}), learning that the difference in the two clusters [\dist{A},\dist{B}] is primarily in the density of \textit{Statewide Wet} realizations.
%
Interestingly, the \textit{Statewide Dry} realizations are both prominently present in both clusters.
%
They visually confirm that the additional occurrences of \textit{Statewide Wet} realizations in one cluster is the driving difference between the two clusters.
%
% The expert also studies the behavior of the two clusters in December, noticing the difference here is in density in the \textit{Statewide Dry} (\req{R2}).
%
Further they learn from the circuit-line diagram that the both the drier and wetter cluster in December evolve to a wetter California in January, but realizes that it primarily evolves from a dry cluster in November considering the number of lines that pass through the respective clusters.
%
Following the dry cluster in November, they learn that this cluster is represented by heavy \textit{Statewide Dry} realizations as well as \textit{Wetter South - Drier North} realizations.
%
Thus the expert concludes that models run that predict a wetter California in January are associated with a overall drier November with a particular dryness in the North of the region, but less clear trends in December.

Another exploration that the expert wishes to undertake is of the behavior of MIROC6, a GCM, in the medium-emission scenario, SSP370, relative to the other GCMs under the same emissions scenario.
%
Exploring through different months, they observe that MIROC6 in February has a distinct pattern compared to the other GCMs in the same emissions scenario .
%
While the other GCMs over February are distributed across \textit{Statewide Wet} and \textit{Statewide Dry} realizations, MIROC6 introduces heavy \textit{Wetter North - Drier South}  and \textit{Wetter South - Drier North} realizations (\req{R1}).
%
This difference in mixture of realizations may warrant further investigation into MIROC6 - SSP370 over February to understand the reasons for the model run's distinct behavior.

\begin{figure}[!t]
    \centering
    % \includegraphics[width=\textwidth]{figures/Dan.pdf}
     \includegraphics[width=\linewidth]{Dan.pdf}
    \caption{The California case study.% with \ClimateSOM. 
     \textit{The Annotated SOM Node Space} used is shown in~\cref{fig:Step1and2}. The expert explored monthly clustering in model runs in January - here the user has hovered over \dist{A}, displaying the lines that pass through it (\req{R3}). The expert explored the cluster behavior by examining its distributions. Lastly, they also explored a single model run, \textit{MIROC6 - SSP370}, against the same SSP for all GCMs (\req{R2}).}
    \label{fig:Dan}
\end{figure}

\paragraph{LOCA-Downscaled CMIP6 over the Northwestern US}
% This section recounts a domain expert's experience using \ClimateSOM to analyze LOCA-downscaled CMIP6 precipitation ensemble over the Northwestern US.
We follow the same process to pick the SOM parameters for the Northwestern US case study, selecting $kR=0.03$ and $kS=0.2$ over a $30\times 30$ SOM grid explaining 78.8\% of the variance.
%
With consultation from an Economist, \textbf{E2}, with experience studying in climate change impacts on the Western Coast of the United States, this case study examines LOCA-downscaled CMIP6 precipitation ensemble over the Northwestern US.
%
We will reference~\cref{fig:Tom} for this case study.
%
First, the expert goes through \step{S1}: Anchor to reorganize the SOM node space to improve the interpretability of the SOM node space and \step{S2}: Annotate and visually identified regions of interest relying previous domain knowledge of precipitation patterns in the Northwestern US.
%
This yielded five annotated regions: (1) \textit{Dry Global}, (2) \textit{Wet NE}, (3) \textit{Dry SW}, (4) \textit{Wet Global}, and (5) Remainder.
% \begin{figure}[h]
%     \centering
%     \includegraphics[width=0.8\linewidth]{figures/TomSpace.png}
%     \caption{The \textit{The Annotated SOM Node Space} for the Northwestern US case study}
%     \label{fig:TomSpace}
% \end{figure}

% The expert first examines the behavior of all GCMs over the historical period in January noticing the breakdown of the selected months into the annotated regions, \textit{Dry Global} and \textit{Remainder} being the most prominent (\req{R1}).
%
% This is seen in \highlight{B}.
%
% Checking the results for the other Fall months, the expert observes the Fall vs.


First, the expert wishes to compare historical and SSP585 precipitation projections in April. 
%
Here, this transition is represented by a change from a unimodal distribution of precipitation realizations to a bimodal distribution (\highlight{C}).
%
In particular, the expert notices a shift to a wetter Northeast (of the region) under this higher emissions future, which they recognize to be a \textit{notable change} with particular significance to agriculture in the region (\req{R2}).
% \begin{figure}[h]
    % \centering
    % \includegraphics[width=\linewidth]{figures/TomBimodal.png}
    % \caption{The shift from a unimodal to bimodal distribution of precipitation realizations in April, comparing historical (left) and SSP585 (right) ensemble model runs.}
    % \label{fig:TomSpace}
% \end{figure}
%
In contrast, in November, the expert observes a movement towards a wetter Northwest (of the region) in a higher emissions future.
%
The expert therefore confirms that the effect on a higher emissions future on precipitation in the Northwestern US is drastically \textbf{different} between April and November, where April expects higher precipitation in the Northeast compared to the Northwest.
%
This is an interesting insight, as simple statistical analysis would not reveal such spatial patterns.


% with visual descriptions of the changes.
% \begin{figure}[h]
%     \centering
%     \includegraphics[width=\linewidth]{figures/TomNov.png}
%     \caption{The shift from a precipitation realizations in November, comparing historical (left) and SSP585 (right) ensemble model runs.}
%     \label{fig:TomSpace}
% \end{figure}
%
% Drilling down to a particular climate model, they observe that for ACCESS-CM2 comparing historical to a higher emission model in November, they see an increase in precpitation in the westside and coastal regions of the Northwestern US (\req{R2}).



Next, the expert wishes to explore whether there are outlier GCMs that exhibit different SSP forcing from historical to SSP585 in November. 
%
Consulting the circuit-line diagram for the GCM clustering of forcings, they observe two clusters of GCMs in November, so they click on the two clusters.
%
Here, they observe that one of the clusters (\highlight{B}), \textit{``homes in the lower right area of the SOM distribution''} thus representing GCMs that predict a drier November with drier anomalies concentrated toward the Southwest of the region.
%
On the other hand, the other cluster of GCMs (\highlight{C}), in contrast, collects toward the lower left area of the SOM distribution, and thus a wetter coastal region and wetter border with Canada, highlighting the two clusters of GCMs in its SSP forcing (\req{R3}).
%
They hypothesize that the the GCM's forcing in the first cluster may be anomalous, and thus may warrant further investigation.
%
Throughout this session, the expert was impressed with \ClimateSOM's ability to \textit{``visualize big picture patterns quickly''} and that it was \textit{``just not something (they) could do otherwise''}.




\begin{figure*}[h]
    \centering
    \includegraphics[width=\textwidth]{Tom-6.pdf}
    \caption{The Northwestern US case study with \ClimateSOM. \highlight{A}: Comparing historical and SSP585 projections in April (\req{R2}), \highlight{B},\highlight{C}: Exploring the two GCM forcing clusters in November and comparing their behavior (\req{R3}).}
    \label{fig:Tom}
\end{figure*}

\subsubsection{Evaluating the LLM integration}\label{ssec:llm_eval}
While past works demonstrate the ability of LLMs to generate SQL queries and a general text summarizer, past works on LLMs have not explicitly evaluated two spatial abilities: (1) the spatial ability to map region identifiers to a list of counties and (2) the spatial ability to summarize a list of counties to a short summary.
%
Additionally, given the possible ambiguity that may arise in how places are defined (i.e., \textit{Southern California} may mean different boundaries to different people), we conducted a short evaluation of our LLM in the above two tasks with 8 participants all familiar with the geography of California.
%
% Specifically, we designed two tasks. 

First, given a common region identifier (e.g., \textit{Southern California}), we asked the participants to evaluate how reasonable the list of counties generated by the LLM was.
%
Second, given a list of counties randomly selected from one of the textual summaries of a SOM node in the California case study, we asked the participants to evaluate how reasonable the regional summaries were.
%
Both tasks were evaluated on a 5-point Likert scale. (\textit{1: Strongly Agree, 2: Agree, 3: Neither Agree nor Disagree, 4: Disagree, 5: Strongly Disagree}).
%
We show examples of the tasks in the Appendix.

Collecting 10 of each tasks (80 responses per task), we found that for T1, the average score was 2.11 (SD=1.32) and for T2, the average score was 2.5 (SD=1.55).
%
This suggests that while the LLMs are not perfect, they are able to generate reasonable lists of counties and summaries of regions that most participants found reasonable.
%
Admittedly, in the first task there were regions where responses were more polarized, especially for vaguely defined region identifiers like \textit{Desert Region of California} and in the second task, when the list of counties were highly disconnected geographically, generating textual summaries of such regions was challenging.
%
Another limitation of the current implementation is that we use counties as the unit of analysis. 
%
However, natural language descriptions of regions do not always conform to county boundaries--some regional identifiers may span or divide multiple counties. As a result, generating finer-grained boundaries and summaries that account for these irregularities is beyond the scope of this integration.

Lastly, we also test the LLM's PostGIS generation for \ClimateSOM by test generating 20 queries for each of the three categories and manually inspecting the results.
%
We found that our LLM's generation was accurate in \textasciitilde 90\% of the cases, with the remaining \textasciitilde 10\% being cases where the LLM generated a query with incorrect syntax.
%
\textbf{In all, we find the failure cases with our LLM integration arising mainly from ambiguous region identifiers for spatial tasks, the county-level unit of analysis, and incorrect syntax generation for the PostGIS queries}.
%
This suggests a conversational interface that can allow the user to correct the LLMs when they generate unreasonable results as a future direction.
